\section{紧合性赋值判据的证明}

于是考虑 $\Cht^{r,d,\ol p\le p}$ 的一个点 $\wt\cE=(\cE\hookrightarrow\cE'\hookleftarrow\cE''\overset{\sim}{\hookleftarrow}{}^\tau\cE)$，它取值于离散赋值环 $A$ 的分式域 $K_A$。要证明：\ul{将 $A$ 换为有限扩张后}，它延拓为 $\Cht^{r,d,\ol p\le p}$ 的唯一的、取值于 $A$ 的点。

这个点必然对应于 $\wt\cE$ 的一般纤维 $V$ 中的一个迭代 $\varphi$-格 $M$；等价地，对应于它在 $\wh V=V\otimes_{K_{A_X}}\wh K_{A_X}$ 中的完备化 $\wh M$。

$V$ 的每个迭代 $\varphi$-格 $M$ 都诱导一个划分 $\ul r_M=(r=r_1+\cdots+r_k)$，以及 $\wh V$ 的一个滤过，由秩为 $s\in\{0,r_1,r_1+r_2,\ldots,r\}=\ul r_M^{-}\cup\ul r_M^{+}$ 的子 $\varphi$-空间 $\wh V_s(M)$ 构成。

映射
\[ \wh W\mapsto\wh W^{\wh M}=\wh M\cap\wh W/\pi_A(\wh M\cap\wh W) \]
\origpage{43}诱导出：与滤过 $(\wh V_s(M))_{s\in\ul r_M^{-}\cup\ul r_M^{+}}$ 相容的（即介于两个相邻的 $\wh V_s(M)$ 之间的）$\wh V$ 的子 $\varphi$-空间 $\wh W$ 的集合，到 $V^M=M/\pi_A M$ 的「好」子空间的集合之间的一个双射。

记 $\cE(M),\cE'(M),\cE''(M)$ 为 $X\otimes A$ 上沿迭代 $\varphi$-格 $M$ 延拓 $\cE,\cE',\cE''$ 所得的丛，$\cE^M,\cE'^M,\cE''^M$ 为它们模 $\pi_A$ 的约化。后者带有递增滤过 $(\cE_s^M),(\cE_s'^M),(\cE_s''^M)$，由秩为 $s\in\ul r_M^{-}\cup\ul r_M^{+}$ 的极大子丛给出；而 ${}^\tau\cE^M$ 带有递减滤过 $(\ol\cE_s^M)$，由秩为 $r-s$ 的极大子丛给出，使得 ${}^\tau\cE_s^M\cap\ol\cE_s^M=0$，$\forall s$。此外，给定 $V^M$ 的一个秩为 $w$ 的「好」子空间，等价于给定 $\cE^M\hookrightarrow\cE'^M\hookleftarrow\cE''^M$ 中一个秩为 $w$ 的「好」子对象 $\cE_w^M\hookrightarrow\cE_w'^M\hookleftarrow\cE_w''^M$。

一个迭代 $\varphi$-格 $M$ 诱导 $\Cht^{r,d,\ol p\le p}(A)$ 上的一点，当且仅当它满足以下条件：

(1) 对每个 $s\in\ul r_M^{-}\cap\ul r_M^{+}$，挠层 ${}^\tau\cE^M/(\ol\cE_s^M\oplus{}^\tau\cE_s^M)$ 的次数为 1。

(2) 对每个 $s\in\ul r_M^{+}$，同态 $\ol\cE_{s^-}^M/\ol\cE_s^M\to\cE_s''^M/\cE_{s^-}''^M$ 处处良定义，且在 $X\otimes\kappa_A$ 上可逆；挠层 $\cE_s'^M/\cE_s^M$ 的次数为 1；并且当 $s<r$ 时 $\cE_s''^M=\cE_s'^M$。

(3) 对每个 $s\in\ul r_M^{-}$，有 ${}^\tau\cE_{s^+}^M+\ol\cE_s^M={}^\tau\cE^M$。

(4) 对每个 $s\in\ul r_M^{-}\cap\ul r_M^{+}$，有 $\deg\cE_s^M=d(s)$。

(5) 对每个秩为 $w$ 的「好」子对象 $\cE_w^M$，有
\[ \deg\cE_w^M\le d(w)\quad\text{若}\quad\cE_w^M\ \text{为 I 型，并且} \]
\[ \deg\cE_w^M\le d(w)-1\quad\text{若}\quad\cE_w^M\ \text{为 II 型。} \]

\begin{lemma*}
对 $V$ 的每个迭代 $\varphi$-格 $M$，有
\begin{itemize}
\item[(i)] 所有挠层 ${}^\tau\cE^M/({}^\tau\cE_s^M\oplus\ol\cE_s^M)$ 的次数为 0 或 1。
\item[(ii)] 若 $s_0\in\ul r_M^{-}\cap\ul r_M^{+}$，则条件 (1) 对所有 $s\ge s_0$ 成立，当且仅当条件 (2) 对所有 $s\ge s_0$ 成立。
\end{itemize}
\end{lemma*}

从一个迭代 $\varphi$-格 $M$ 过渡到另一个，靠的是此处与 $V^M$ 的「好」子空间 $W$ 相联系的、「Langton 式」的初等变换。

\origpage{44}\begin{proposition*}
设 $W$ 是 $V^M$ 的一个秩为 $w_0$ 的「好」子空间（或者等价地，是 $\cE^M$ 的一个「好」子对象 $\cE_{w_0}^M$，或是 $\wh V$ 的一个「好」子 $\varphi$-空间 $\wh W$），并设
\[ M'=\mathrm{Ker}\,[M\to M/\pi_A M=V^M\to V^M/W] \]
是 $M$ 经 $W$ 或 $\wh W$ 的（初等）变换。那么
\begin{itemize}
\item[(i)] 有典范同构
\[ \ol\cE_{w_0}^{M'}\cong{}^\tau\cE^M/{}^\tau\cE_{w_0}^M,\qquad{}^\tau\cE_{w_0}^M\cong{}^\tau\cE^{M'}/\ol\cE_{w_0}^{M'}. \]
\item[(ii)] 有典范嵌入
\[ \cE_{w_0}^{M'}\hookrightarrow\cE_{w_0}^M,\qquad\cE^M/\cE_{w_0}^M\hookrightarrow\cE^{M'}/\cE_{w_0}^{M'}, \]
其余核为挠，且次数相同，等于 0 或 1。该次数等于 1，当且仅当
\[ {}^\tau\cE_{w_0}^{M'}\oplus\ol\cE_{w_0}^{M'}\subsetneq{}^\tau\cE^{M'}. \]
\end{itemize}
\end{proposition*}

在这种情形下，与 $M'$ 相伴的 $r$ 的划分与 $\wh V$ 的滤过由下式给出：
\[ \ul r_{M'}^{+}=\ul r_M^{+}\cup\{w_0\}\qquad\{\wh V_{s'}(M')\}=\{\wh V_s(M)\}\cup\{\wh W\}. \]

称 $M$ 是 $M'$ 的一个互反变换。它由 $M'$ 与秩 $w_0\in\ul r_{M'}^{-}\cap\ul r_{M'}^{+}$ 唯一确定。

还考虑有限的迭代 $\varphi$-格链
\[ M=M_0\supsetneq M_1\supsetneq\cdots\supsetneq M_k=M', \]
其中每一个都是前一个经同一个秩为 $w_0$ 的 $\wh V$ 的子 $\varphi$-空间 $\wh W$ 所作的初等变换。若
\[ \cE_{w_0}^{M_0}=\cE_{w_0}^{M_1}=\cdots=\cE_{w_0}^{M_{k-1}}\supsetneq\cE_{w_0}^{M_k}, \]
则称 $M'$ 是 $M$ 经 $\wh W$ 的严格变换；若
\[ \cE_{w_0}^{M_0}=\cE_{w_0}^{M_1}=\cdots=\cE_{w_0}^{M_k}, \]
且下列两个条件之一成立：
\begin{itemize}
\item \origpage{45}或者 $\wh W$ 不出现在与 $M$ 相伴的 $\wh V$ 的典范滤过 $\{\wh V_s(M)\}_s$ 中，
\item 或者
\[ {}^\tau\cE_{w_0}^M\oplus\ol\cE_{w_0}^M\subsetneq{}^\tau\cE^M, \]
\end{itemize}
则称 $M$ 是 $M'$ 经 $\wh W$（或在秩 $w_0$ 处）的极大修改。

最后，若存在 $k'$，$0<k'\le k$，使得
\[ \cE_{w_0}^{M_0}=\cdots=\cE_{w_0}^{M_{k'-1}}\supsetneq\cE_{w_0}^{M_{k'}}=\cdots=\cE_{w_0}^{M_k}, \]
且 $M_0$ 是 $M_{k'-1}$ 在秩 $w_0$ 处的极大修改，则称 $M$ 是 $M'$ 经 $\wh W$（或在 $w_0$ 处）的严格互反变换。

下面两个基本结果保证严格变换、极大修改与严格互反变换的存在性：

\begin{proposition*}
设 $M$ 是 $V$ 的一个迭代 $\varphi$-格，$\wh W$ 是 $\wh V$ 的一个秩为 $w_0$、与滤过 $\{\wh V_s(M)\}$ 相容的子 $\varphi$-空间，并且
\[ \deg\cE_{w_0}^M>d(w_0). \]
那么 $M$ 具有一个经 $\wh W$ 的严格变换。
\end{proposition*}

\begin{proof}
这是 \cite{lafforgue98} 第 2d 段的命题 14。
\end{proof}

\begin{proposition*}
设 $\{\wh W_w\}$ 是 $\varphi$-空间 $V$ 的一个滤过。

那么，对每个相伴的 $\wh V$ 滤过包含于 $\{\wh W_w\}$ 的迭代 $\varphi$-格 $M$，以及对每个元素 $\wh W_{w_0}\in\{\wh W_w\}$，$M$ 或者具有一个严格互反变换（从而具有一个相伴滤过包含 $\wh W_{w_0}$ 的极大修改），或者具有一个相伴滤过不含 $\wh W_{w_0}$ 的极大修改；这里只需将 $A$ 换为一个有限扩张（可能分歧），它只依赖于所有子商 $\cE_{w^{+}}^M/\cE_w^M$ 的极小斜率 $\mu^{-}$ 与极大斜率 $\mu^{+}$。
\end{proposition*}

\begin{proof}
若 $\wh W_{w_0}$ 出现在与 $M$ 相伴的 $\wh V$ 滤过中，则 $M$ 经 $\wh W_{w_0}$ 的互反变换之存在性，就是 \cite{lafforgue98} 第 2d 段的命题 15，并按 \cite{lafforgue02} 命题 III.11 之后的注记作了修正。证明的过程是：在诸 $\cE^M$ 上构造水平结构 $I=N$（其次数足够大，依赖于所有 $\mu^{+}$ 与 $\mu^{-}$）。我们选取水平 $I=N$，使其位于极点、零点之外\footnote{译注：原文此处跨页，句子未完，接下页。}。
\end{proof}
